\section{Conclusion}\label{sec:conclusion}
We asked why accurate multimodal segmentation models fail when one sensor degrades. Controlled comparisons on DELIVER showed that their heavy dependence on depth is learned from clean training data: it appears in our model without depth supervision as strongly as in DGFusion with it, and it largely disappears in both architectures when degraded inputs are added to training. Built on this finding, \ours{} combines a frozen DINOv3 backbone, symmetric fusion, and a high-resolution path with two-pass training, and improves clean accuracy and robustness together on DELIVER \todo{and on MUSES and MCubeS, once measured}. The analysis also exposes what remains: the high-resolution path recovers thin objects but lets fine detail mislead the per-pixel classifier, as in posters read as traffic lights, which suggests routing fine detail to mask prediction rather than to per-pixel classification. Measuring the same dependence on real depth sensors is the next step that our synthetic-data findings call for.
